\documentclass[11pt,a4paper]{article}
% preamble.tex — estilo común de todos los apuntes Froth.
% Cada apunte hace % preamble.tex — estilo común de todos los apuntes Froth.
% Cada apunte hace % preamble.tex — estilo común de todos los apuntes Froth.
% Cada apunte hace \input{preamble} justo después de \documentclass.
\usepackage[margin=2.4cm]{geometry}
\usepackage{xcolor}
\definecolor{litblue}{RGB}{31,79,121}
\definecolor{litgreen}{RGB}{27,94,32}
\definecolor{litorange}{RGB}{230,81,0}
\usepackage{parskip}
\usepackage{titlesec}
\titleformat{\section}{\Large\bfseries\color{litblue}}{\thesection}{0.6em}{}
\titleformat{\subsection}{\large\bfseries\color{litblue}}{\thesubsection}{0.6em}{}
\usepackage{tcolorbox}
\usepackage{fancyhdr}
\pagestyle{fancy}
\fancyhf{}
\fancyhead[L]{\small\color{litblue}Froth — Apuntes de ML}
\fancyhead[R]{\small\thepage}
\renewcommand{\headrulewidth}{0.4pt}
\usepackage[colorlinks=true,linkcolor=litblue,urlcolor=litblue]{hyperref}

% Cabecera de cada apunte: \cabecera{numero}{titulo}
\newcommand{\cabecera}[2]{%
  \begin{tcolorbox}[colback=litblue,coltext=white,colframe=litblue,arc=2mm]
    {\Large\bfseries Apunte #1 — #2}\\[3pt]
    {\small Proyecto Froth · aprender Machine Learning construyendo}
  \end{tcolorbox}\vspace{0.8em}}

% Cajas temáticas
\newtcolorbox{concepto}[1]{colback=blue!4,colframe=litblue,title={Concepto: #1},arc=1.5mm}
\newtcolorbox{practica}{colback=green!5,colframe=litgreen,title={Pruébalo tú (teclado en mano)},arc=1.5mm}
\newtcolorbox{ojo}{colback=orange!8,colframe=litorange,title={Ojo / error típico},arc=1.5mm}
\newtcolorbox{resumen}{colback=gray!8,colframe=gray!60,title={En una frase},arc=1.5mm}

\newcommand{\codigo}[1]{\texttt{#1}}
 justo después de \documentclass.
\usepackage[margin=2.4cm]{geometry}
\usepackage{xcolor}
\definecolor{litblue}{RGB}{31,79,121}
\definecolor{litgreen}{RGB}{27,94,32}
\definecolor{litorange}{RGB}{230,81,0}
\usepackage{parskip}
\usepackage{titlesec}
\titleformat{\section}{\Large\bfseries\color{litblue}}{\thesection}{0.6em}{}
\titleformat{\subsection}{\large\bfseries\color{litblue}}{\thesubsection}{0.6em}{}
\usepackage{tcolorbox}
\usepackage{fancyhdr}
\pagestyle{fancy}
\fancyhf{}
\fancyhead[L]{\small\color{litblue}Froth — Apuntes de ML}
\fancyhead[R]{\small\thepage}
\renewcommand{\headrulewidth}{0.4pt}
\usepackage[colorlinks=true,linkcolor=litblue,urlcolor=litblue]{hyperref}

% Cabecera de cada apunte: \cabecera{numero}{titulo}
\newcommand{\cabecera}[2]{%
  \begin{tcolorbox}[colback=litblue,coltext=white,colframe=litblue,arc=2mm]
    {\Large\bfseries Apunte #1 — #2}\\[3pt]
    {\small Proyecto Froth · aprender Machine Learning construyendo}
  \end{tcolorbox}\vspace{0.8em}}

% Cajas temáticas
\newtcolorbox{concepto}[1]{colback=blue!4,colframe=litblue,title={Concepto: #1},arc=1.5mm}
\newtcolorbox{practica}{colback=green!5,colframe=litgreen,title={Pruébalo tú (teclado en mano)},arc=1.5mm}
\newtcolorbox{ojo}{colback=orange!8,colframe=litorange,title={Ojo / error típico},arc=1.5mm}
\newtcolorbox{resumen}{colback=gray!8,colframe=gray!60,title={En una frase},arc=1.5mm}

\newcommand{\codigo}[1]{\texttt{#1}}
 justo después de \documentclass.
\usepackage[margin=2.4cm]{geometry}
\usepackage{xcolor}
\definecolor{litblue}{RGB}{31,79,121}
\definecolor{litgreen}{RGB}{27,94,32}
\definecolor{litorange}{RGB}{230,81,0}
\usepackage{parskip}
\usepackage{titlesec}
\titleformat{\section}{\Large\bfseries\color{litblue}}{\thesection}{0.6em}{}
\titleformat{\subsection}{\large\bfseries\color{litblue}}{\thesubsection}{0.6em}{}
\usepackage{tcolorbox}
\usepackage{fancyhdr}
\pagestyle{fancy}
\fancyhf{}
\fancyhead[L]{\small\color{litblue}Froth — Apuntes de ML}
\fancyhead[R]{\small\thepage}
\renewcommand{\headrulewidth}{0.4pt}
\usepackage[colorlinks=true,linkcolor=litblue,urlcolor=litblue]{hyperref}

% Cabecera de cada apunte: \cabecera{numero}{titulo}
\newcommand{\cabecera}[2]{%
  \begin{tcolorbox}[colback=litblue,coltext=white,colframe=litblue,arc=2mm]
    {\Large\bfseries Apunte #1 — #2}\\[3pt]
    {\small Proyecto Froth · aprender Machine Learning construyendo}
  \end{tcolorbox}\vspace{0.8em}}

% Cajas temáticas
\newtcolorbox{concepto}[1]{colback=blue!4,colframe=litblue,title={Concepto: #1},arc=1.5mm}
\newtcolorbox{practica}{colback=green!5,colframe=litgreen,title={Pruébalo tú (teclado en mano)},arc=1.5mm}
\newtcolorbox{ojo}{colback=orange!8,colframe=litorange,title={Ojo / error típico},arc=1.5mm}
\newtcolorbox{resumen}{colback=gray!8,colframe=gray!60,title={En una frase},arc=1.5mm}

\newcommand{\codigo}[1]{\texttt{#1}}

\begin{document}
\cabecera{29}{Fase 6: generalizar no es entrenar (y el motor multi-topic)}

\section{La aclaración que abre la fase}

El plan de Batman: ``iterar cada uno de los 22 títulos del cohorte y seguir entrenando''.
Correcto en espíritu — pero son \textbf{dos actos distintos}:

\begin{concepto}{generalizar (6.1--6.3) vs.\ entrenar (6.5)}
\begin{itemize}
  \item \textbf{Iterar los 22 títulos NO entrena nada}: el modelo (SPECTER) no cambia ni un
        peso. Lo que se endurece es el \textbf{código} — cada título raro que rompe algo
        obliga a arreglarlo. Es la escalera 1$\rightarrow$2$\rightarrow$22 del plan original.
  \item \textbf{Entrenar de verdad} (6.5) sí cambia los pesos del modelo: fine-tuning de
        SPECTER con los miles de abstracts acumulados de los 22 corpus. PyTorch, loss,
        épocas — en la GPU de Batman (RTX 5060, confirmada capaz).
\end{itemize}
Los dos actos se alimentan: la generalización produce el dataset del entrenamiento.
\end{concepto}

\section{El motor multi-topic (6.1)}

Hasta ayer el topic vivía \emph{clavado} en \codigo{config.py}. Ahora:

\begin{concepto}{config.set\_topic(título)}
Una llamada re-apunta TODO el motor: deriva términos de búsqueda nuevos y mueve las rutas de
datos a una carpeta propia (\codigo{2\_Datos/topics/<slug>/}). Funciona porque cada módulo
lee \codigo{config.*} \emph{en el momento de ejecutar}, no al importar — un detalle de
diseño que pagó dividendos. (Trampa clásica cazada: \codigo{def f(x=config.TOPIC)} congela
el valor al importar; se corrigió a \codigo{x=None} + resolver adentro.)
\end{concepto}

\begin{concepto}{derive\_search\_terms(): de título a búsquedas}
Los términos del topic 1 se escribieron a mano. Para CUALQUIER título, la heurística de
\codigo{froth/terms.py}: quita palabras de relleno (``innovative approaches for the...''),
conserva \textbf{frases reales del título} (bigramas/trigramas consecutivos: \emph{froth
flotation}, \emph{scheelite ores}, \emph{Quantitative XRD}), respeta \textbf{acrónimos}
(XRD, LIBS, NiMH) y extrae el \textbf{contenido de paréntesis} (ahí viven los elementos y
yacimientos). Determinista, sin dependencias. La capa LLM local (Ollama) vendrá encima como
refinador opcional — decisión ya tomada: local, privado, gratis.
\end{concepto}

\section{El hito 1\texorpdfstring{$\rightarrow$}{->}2 (probado)}

Pipeline completo sobre \emph{``scheelite ores by froth flotation''} \textbf{sin tocar
código}: 47 papers, carpeta propia, y clusters con sentido metalúrgico (concentrador
Falcon, reactivos SNF, espumantes). Y la prueba pagó de inmediato:

\begin{ojo}
\textbf{Lo que el 2º topic enseñó} (por esto existe la escalera):
\begin{enumerate}
  \item La keyword ``nbsp'' delató que los abstracts de Crossref traen \textbf{entidades
        HTML} ($\&$nbsp;) — el limpiador ahora hace \codigo{html.unescape}.
  \item El \codigo{.gitignore} no cubría las carpetas nuevas por topic — los datos casi
        acaban en GitHub.
  \item En corpus chicos (47 papers) la granularidad fija de 8 deja demasiado ruido
        $\rightarrow$ el pipeline ahora acepta \codigo{min\_cluster\_size="auto"} (elige
        por DBCV, Apunte 27, por corpus).
\end{enumerate}
Tres arreglos reales con UN solo título nuevo. Los 22 del batch encontrarán más — ese es
exactamente su trabajo.
\end{ojo}

\section{El desfile (6.2)}

\codigo{froth/batch.py} lee los 22 títulos del CSV del cohorte y corre el pipeline para
cada uno, con reporte de robustez por topic (papers, relevantes, clusters, ruido,
granularidad elegida, segundos, estado). Diseño defensivo: el reporte se guarda tras
\textbf{cada} topic y las corridas siguientes \textbf{reanudan} donde quedaron (un
rate-limit no pierde nada). Los fallos no detienen el desfile: quedan como FAILED en la
tabla — y esa tabla es la agenda de arreglos de la 6.3.

\begin{resumen}
Iterar títulos endurece el código (generalizar); cambiar pesos es entrenar (6.5). El motor
ya es multi-topic (set\_topic + términos derivados + carpetas por topic), probado con
scheelite ores — que de paso regaló 3 arreglos. El batch de 22 corre con reanudación y
reporte por topic: su tabla de fallos es el plan de trabajo de la 6.3.
\end{resumen}

\end{document}
